\documentclass[aspectratio=169,11pt]{beamer}

\usepackage{fontspec}
\usepackage{booktabs}
\usepackage{array}
\usepackage{ragged2e}
\usepackage{hyperref}

\setsansfont{Helvetica Neue}
\setmonofont{Menlo}
\usetheme{Boadilla}
\usecolortheme{dolphin}
\setbeamertemplate{navigation symbols}{}
\setbeamertemplate{itemize item}{\raise0.2ex\hbox{\tiny$\blacktriangleright$}}
\setbeamercolor{structure}{fg=blue!55!black}
\setbeamercolor{block title}{fg=white,bg=blue!55!black}
\setbeamercolor{block body}{bg=blue!4}
\setbeamerfont{frametitle}{series=\bfseries}
\setbeamerfont{title}{series=\bfseries}

\title{MPPI Local Planner trên Companion Computer}
\subtitle{Tích hợp LiDAR 3D với ArduPilot GUIDED qua MAVLink}
\author{Nguyễn Thành Lâm}
\institute{Báo cáo tiến độ kỹ thuật}
\date{09/09/2026}

\newcommand{\ok}{\textcolor{green!45!black}{\textbf{Hoàn thành}}}
\newcommand{\partdone}{\textcolor{orange!80!black}{\textbf{Một phần}}}
\newcommand{\todo}{\textcolor{red!65!black}{\textbf{Chưa làm}}}

\begin{document}

\begin{frame}
  \titlepage
\end{frame}

\begin{frame}{Mục tiêu và quyết định kiến trúc}
  \begin{block}{Vấn đề}
    Pipeline \texttt{AP\_Avoidance} yêu cầu tiền xử lý obstacle theo giao thức
    của ArduPilot. Cách này làm perception và tracking phụ thuộc vào autopilot.
  \end{block}

  \vspace{0.4em}
  \textbf{Thiết kế đã chọn}
  \begin{itemize}
    \item Companion computer nhận point cloud LiDAR 3D và chạy local planner.
    \item MPPI sinh $u=[v_x,v_y,v_z,\dot\psi]$ theo thời gian thực.
    \item ArduPilot ở \texttt{GUIDED} chỉ bám velocity setpoint.
    \item Avoidance nội bộ của ArduPilot được tắt hoàn toàn.
  \end{itemize}
\end{frame}

\begin{frame}{Luồng dữ liệu đã triển khai}
  \centering
  \begin{minipage}{0.92\textwidth}
    \centering
    \texttt{LiDAR 3D PointCloud}\quad +\quad
    \texttt{Odometry / MAVLink telemetry}

    \vspace{0.7em}
    $\Downarrow$

    \vspace{0.7em}
    \fbox{\parbox{0.78\textwidth}{\centering
      Downsample và đổi frame FLU--FRD--ENU/NED\\[0.35em]
      MPPI: 7-state, 4-action, 500 samples, horizon 3 s}}

    \vspace{0.7em}
    $\Downarrow$

    \vspace{0.7em}
    \texttt{SET\_POSITION\_TARGET\_LOCAL\_NED}\quad
    \texttt{type\_mask = 1479}

    \vspace{0.7em}
    $\Downarrow$

    \vspace{0.7em}
    \textbf{ArduPilot GUIDED: attitude control và velocity tracking}
  \end{minipage}
\end{frame}

\begin{frame}{Kết quả thực hiện}
  \renewcommand{\arraystretch}{1.22}
  \begin{tabular}{>{\RaggedRight\arraybackslash}p{0.38\textwidth}
                  >{\RaggedRight\arraybackslash}p{0.22\textwidth}
                  >{\RaggedRight\arraybackslash}p{0.30\textwidth}}
    \toprule
    Hạng mục & Trạng thái & Bằng chứng chính \\
    \midrule
    Tách avoidance khỏi ArduPilot & \ok & \texttt{OA\_TYPE=0}, \texttt{AVOID\_ENABLE=0} \\
    Vanilla MPPI local planner & \ok & state 7D, action 4D, predicted path \\
    MAVLink velocity + yaw rate & \ok & mask 1479, gửi 10 Hz \\
    LiDAR 3D trong Gazebo & \ok & lọc, downsample, đổi frame \\
    Global reference & \partdone & waypoint nhập tay \\
    Edge computer / quad thật & \partdone & có state MAVLink, chưa HIL/flight test \\
    PA-MPPI & \partdone & v0 + rigid-body experimental \\
    \bottomrule
  \end{tabular}
\end{frame}

\begin{frame}{Bằng chứng kiểm thử}
  \begin{columns}[T,onlytextwidth]
    \begin{column}{0.49\textwidth}
      \begin{block}{Offline MPPI, chạy lại 09/09}
        \begin{itemize}
          \item Tới đích: \textbf{Có}
          \item Vanilla: sai số \textbf{0.91 m}, clearance \textbf{6.09 m}
          \item PA-MPPI v0: sai số \textbf{0.86 m}
          \item PA-MPPI v0: clearance \textbf{5.85 m}
          \item Cả hai đạt ngưỡng \textbf{5.00 m}
        \end{itemize}
      \end{block}
    \end{column}
    \begin{column}{0.49\textwidth}
      \begin{block}{Tích hợp}
        \begin{itemize}
          \item Gazebo plugin build thành công
          \item Mask khớp enum \texttt{pymavlink}
          \item \textbf{18/18} automated tests đạt
          \item SITL arm, takeoff 20 m, hold, land đạt
        \end{itemize}
      \end{block}
    \end{column}
  \end{columns}

  \vspace{0.5em}
  \alert{Lưu ý:} log bay MPPI cũ nằm trong \texttt{/tmp}. Cần quay và lưu lại
  video cùng log trong \texttt{output/} để tạo bằng chứng tái lập.
\end{frame}

\begin{frame}{Đánh giá kỹ thuật}
  \begin{columns}[T,onlytextwidth]
    \begin{column}{0.49\textwidth}
      \textbf{Điểm mạnh}
      \begin{itemize}
        \item Phân tầng đúng giữa planner và autopilot
        \item Sửa đúng lỗi mask làm AP bỏ qua velocity
        \item Module hóa controller, sensor và MAVLink
        \item Có hold khi sensor stale hoặc obstacle quá gần
      \end{itemize}
    \end{column}
    \begin{column}{0.49\textwidth}
      \textbf{Rủi ro còn lại}
      \begin{itemize}
        \item Nhánh MAVLink mới xoay LiDAR theo yaw
        \item Rigid branch chưa bay hover gate trong Gazebo
        \item Mapper gần-field còn là giả thiết LiDAR 360 độ
        \item Chưa benchmark trên edge computer
        \item Waypoint tay chưa phải global planner
      \end{itemize}
    \end{column}
  \end{columns}

  \vspace{0.6em}
  \begin{block}{Mức hoàn thành}
    \textbf{8.5/10 cho prototype SITL}; khoảng \textbf{60\% cho triển khai trên
    quad thật}.
  \end{block}
\end{frame}

\begin{frame}{Kế hoạch vòng tiếp theo}
  \begin{enumerate}
    \item Quay demo Gazebo + RViz Path, sampled rollouts và MAVLink target.
    \item So sánh vanilla với PA-MPPI v0 trên cùng kịch bản và random seed.
    \item Hoàn thiện attitude quaternion cho nhánh edge computer.
    \item Benchmark 10 Hz: mean, p95, worst-case và
      deadline miss.
    \item Tích hợp hit/miss ray từ sensor thật và kiểm thử dropout.
    \item Port thrust/body-rate dynamics nếu cần tái lập đầy đủ paper.
  \end{enumerate}

  \vfill
  \footnotesize
  Mã nguồn: \url{https://github.com/thanhlamauto/ardu_plugin_gazebo}
\end{frame}

\begin{frame}{Phạm vi PA-MPPI v0}
  \small
  \renewcommand{\arraystretch}{1.15}
  \begin{tabular}{>{\RaggedRight\arraybackslash}p{0.28\textwidth}
                  >{\RaggedRight\arraybackslash}p{0.30\textwidth}
                  >{\RaggedRight\arraybackslash}p{0.32\textwidth}}
    \toprule
    Thành phần & Port v0 trong repo & Bài báo PA-MPPI \\
    \midrule
    Bản đồ & voxel chưa biết/tự do/vật cản & ROG-Map, ray casting \\
    Perception cost & hướng goal, thưởng vùng chưa biết, phạt vật cản
      & hai pha theo line-of-sight \\
    Dynamics & v0 $[p,v,\psi]$; rigid $[p,q,v,\omega]$
      & rigid-body, thrust và body rates \\
    Tối ưu & PyTorch, 512 rollout, target 50 Hz
      & JAX, 17,500 rollout, 50 Hz control \\
    Autopilot & \texttt{SET\_ATTITUDE\_TARGET}, chưa bay live
      & low-level controller riêng \\
    \bottomrule
  \end{tabular}

  \vspace{0.35em}
  \alert{Kết luận:} rigid rollout và MAVLink output đã có; cần qua hover gate
  Gazebo trước khi thử tránh vật cản.

  \vfill
  \scriptsize Zhai, Reiter, Scaramuzza, ``Perception-Aware MPPI for Quadrotor
  Exploration,'' IEEE RA-L, 2026.
\end{frame}

\end{document}
